\chapter{Core Units of Analysis}
\label{ch:core-units}

\epigraph{All models are wrong, but some are useful.}{George E.\ P.\ Box}

\noindent
Doxonomics requires a precise vocabulary of parts.
Before we can model how beliefs are produced, we need formal definitions of the entities involved: who produces belief, through what institutional machinery, via which channels, with what resources, carrying what content, reaching which audiences, and feeding back into what structures.

This chapter defines seven core units (agents, institutions, channels, budgets, narratives, audiences, and feedback loops) and gives each a mathematical specification.
Together they constitute the component library from which doxonomic systems are assembled.


\section{Agents}
\label{sec:ch4-agents}

\begin{definition}[Doxonomic Agent]
\label{def:agent}
\index{agent}
A \emph{doxonomic agent} is an entity $a \in \agentspace$ characterized by:
\begin{itemize}[nosep]
  \item a \emph{type} $\tau(a) \in \{\text{funder}, \text{institution}, \text{intermediary}, \text{audience member}\}$,
  \item a \emph{resource endowment} $r(a) \in \R_{\geq 0}$,
  \item a \emph{belief state} $\pi_a \in \Delta(\Omega)$ (a probability distribution over propositions),
  \item a \emph{strategic objective} $u_a: \beliefspace \to \R$ (a utility function over population-level belief states).
\end{itemize}
\end{definition}

The four types play structurally different roles:

\textbf{Funders} allocate resources to channels.
They have preferences about what the population believes, preferences that may be ideological (a foundation promoting free-market ideas), commercial (a corporation promoting consumerism), or governmental (a state promoting national identity).
Their strategic problem is: given budget $r(a)$, how to allocate across channels to maximize $u_a(\beliefvec)$.

\textbf{Institutions} process funding into narrative outputs.
They are the factories of belief production: universities, think tanks, media organizations, PR firms, religious organizations, government agencies.
Their output depends not only on funding but on their internal norms, personnel, accumulated credibility, and organizational culture.

\textbf{Intermediaries} amplify, filter, translate, or repackage narratives as they move through the system.
A journalist who writes up a think-tank report, a social media influencer who popularizes an academic finding, a textbook author who translates research into pedagogical material: each is an intermediary who transforms the narrative in the process of transmitting it.

\textbf{Audience members} receive, interpret, and may resist or adopt narratives.
They are not passive.
They bring prior beliefs, identity commitments, lived experience, and cognitive resources to the encounter with any narrative.
The same narrative may persuade one audience segment and provoke backlash in another.

\begin{remark}
\label{rem:agent-dual-roles}
In practice, agents often occupy multiple roles simultaneously.
A university professor is both an institution (producing research) and an intermediary (translating it for students) and an audience member (consuming media).
The type classification is analytical, not existential: it describes the \emph{role} an agent plays in a specific doxonomic transaction, not a permanent identity.
\end{remark}

The key structural asymmetry is that funders optimize over $\beliefspace$ (they have preferences about what the population believes and resources to act on those preferences), while audience members are typically modeled as updating agents without strategic control over the information environment.
This asymmetry is the source of doxonomic power.

\begin{example}[Agent Asymmetry in Climate Discourse]
\label{ex:climate-agents}
In the climate-belief system:
\begin{itemize}[nosep]
  \item \textbf{Funders}: fossil fuel companies (ExxonMobil, Koch Industries) invested hundreds of millions in climate-doubt production \autocite{oreskes2010}.
  \item \textbf{Institutions}: think tanks (Heartland Institute, Global Warming Policy Foundation) processed funding into contrarian research, op-eds, and congressional testimony.
  \item \textbf{Intermediaries}: cable news producers, op-ed editors, and social media algorithms amplified doubt narratives.
  \item \textbf{Audience}: the general public, which by 2010 was significantly more skeptical about climate change than the scientific consensus warranted.
\end{itemize}
The funding asymmetry: the fossil fuel industry's doxonomic budget for climate doubt dwarfed the combined science-communication budgets of all climate research institutions.
This did not make the doubt narrative \emph{true}; it made it \emph{ubiquitous}.
\end{example}


\section{Institutions as Belief-Processing Functions}
\label{sec:ch4-institutions}

Institutions are the central machinery of doxonomic production.
They are not mere conduits; they \emph{transform} inputs.
A think tank does not simply relay its donor's wishes.
It translates them into credentialed, publishable, quotable form, adding the epistemic capital that makes the output persuasive.

\begin{definition}[Institutional Processing Function]
\label{def:institution-function}
\index{institution!processing function}
An institution $\iota \in \insteco$ is a function
\[
  \iota: \R_{\geq 0} \times \narspace_{\text{in}} \to \narspace_{\text{out}} \times [0,1]
\]
mapping a funding level and input narrative material to an output narrative and a credibility multiplier.
The credibility multiplier reflects the institution's accumulated epistemic capital $\ecap(\iota)$.
\end{definition}

The processing function captures three things:
\begin{enumerate}[nosep]
  \item \textbf{Content transformation}: the input (a donor's policy preference) becomes output (a research paper, a textbook chapter, an op-ed). The content may be identical, but the form is different, and form determines doxonomic force.
  \item \textbf{Credibility attachment}: the institutional brand is affixed to the output. ``A study by the Heritage Foundation'' and ``a study by Harvard'' carry different credibility weights, even if the methodology is identical.
  \item \textbf{Audience routing}: different institutions have access to different distribution networks. A think tank's output reaches policymakers and journalists; a university's output reaches students and academics; an advertising agency's output reaches consumers.
\end{enumerate}

\begin{example}[Credibility Differential]
\label{ex:think-tank-credibility}
A think tank $\iota$ receives funding $\funding{k} = \$2\text{M/year}$ and raw policy positions from its donors.
It outputs research papers, op-eds, and expert testimony.
Its credibility multiplier $\credibility{\iota} = 0.7$ reflects its reputation among journalists and policymakers.
A university economics department performing similar work might have $\credibility{\iota'} = 0.9$.
The difference in $\credibility{}$ means the same proposition, backed by the same evidence, has different doxonomic force depending on its institutional origin.

This credibility differential is what makes \emph{credibility laundering}\index{credibility laundering} possible: a low-credibility funder channels money through a high-credibility institution, gaining the credibility premium without producing the underlying quality.
The institution's brand is the product being purchased.
\end{example}

\begin{definition}[Institutional Capacity]
\label{def:capacity}
\index{institutional capacity}
The \emph{capacity} of an institution $\iota$ is the maximum narrative throughput it can produce per unit time:
\[
  \Theta^{\max}(\iota) = \max_{\funding{}} \left\{ \reach{\iota}(\funding{}) \cdot \credibility{\iota} \cdot \persistence{\iota} \right\}
\]
Capacity depends on staffing, infrastructure, accumulated expertise, and institutional reputation.
It cannot be created instantly: building a credible institution takes years or decades, which is why doxonomic systems exhibit strong path dependence.
\end{definition}

\begin{remark}[Institutional Autonomy]
\label{rem:autonomy}
Institutions are not perfectly obedient to their funders.
A university department funded by a corporation may produce research that contradicts the funder's preferences.
A media outlet owned by a billionaire may employ journalists who resist editorial direction.
This relative autonomy is an important feature, not a bug, of the model: it means the production function $\iota$ has its own internal logic that may partially resist the input signal.
The degree of institutional autonomy varies: a PR firm has near-zero autonomy (it produces what the client wants); a tenured university department has considerable autonomy (funding affects what is studied, not necessarily what is concluded).
\end{remark}


\section{Channels}
\label{sec:ch4-channels}

Channels are the transmission media through which institutional outputs reach audiences.
Each channel has characteristic properties that determine its doxonomic profile.

\begin{definition}[Doxonomic Channel]
\label{def:channel}
\index{channel}
A \emph{channel} is a medium of narrative transmission characterized by:
\begin{itemize}[nosep]
  \item \emph{reach} $\reach{k} \in [0,1]$: the fraction of the population exposed per unit time,
  \item \emph{credibility} $\credibility{k} \in [0,1]$: the weight audiences assign to messages from this channel,
  \item \emph{persistence} $\persistence{k} \in \R_{>0}$: the half-life of the narrative in audience memory,
  \item \emph{cost function} $c_k: \R_{\geq 0} \to [0,1]$: the mapping from funding to effective reach (typically concave: diminishing returns to spending).
\end{itemize}
\end{definition}

Different channels occupy different regions of the reach-credibility-persistence space:

\begin{center}
\begin{tikzpicture}[
  every node/.style={font=\small},
  ch/.style={draw, thick, rounded corners, minimum width=2.8cm, minimum height=0.7cm, align=center, fill=doxoblue!8},
]
\matrix[row sep=0.4cm, column sep=0.3cm] {
  \node[ch] {Education}; &
  \node[font=\small, text width=8.5cm] {Low reach/time, high credibility ($\sim0.85$), very high persistence ($\sim$decades). Shapes foundational priors. Most expensive per-person but most durable.}; \\
  \node[ch] {Prestige Media}; &
  \node[font=\small, text width=8.5cm] {Medium reach, high credibility ($\sim0.7$), medium persistence ($\sim$months). Sets the daily agenda for elites and opinion leaders.}; \\
  \node[ch] {Broadcast Media}; &
  \node[font=\small, text width=8.5cm] {High reach, medium credibility ($\sim0.5$), low persistence ($\sim$days). Mass exposure with high turnover.}; \\
  \node[ch] {Social Media}; &
  \node[font=\small, text width=8.5cm] {Very high reach, low credibility ($\sim0.2$), very low persistence ($\sim$hours). Amplifies, fragments, and accelerates.}; \\
  \node[ch] {Think Tanks}; &
  \node[font=\small, text width=8.5cm] {Low direct reach, high credibility in policy circles ($\sim0.75$), medium persistence. Produces expert legitimacy for policy narratives.}; \\
  \node[ch] {Entertainment}; &
  \node[font=\small, text width=8.5cm] {High reach, low perceived persuasive intent, medium persistence. Operates at layers 3--5: normalizes through representation and affect.}; \\
  \node[ch] {Peer Networks}; &
  \node[font=\small, text width=8.5cm] {Variable reach (local), high credibility ($\sim0.8$), high persistence. Most resistant to top-down manipulation but most responsive to social proof.}; \\
};
\end{tikzpicture}
\end{center}

\begin{remark}[The Credibility-Reach Tradeoff]
\label{rem:tradeoff}
There is a general tradeoff between credibility and reach.
Channels with high credibility (education, peer networks) tend to have limited reach per unit time.
Channels with high reach (broadcast, social media) tend to have lower credibility.
This tradeoff shapes doxonomic strategy: a funder seeking deep belief change (common-sense capture) must invest in high-credibility, high-persistence channels; a funder seeking temporary opinion shifts (e.g., before an election) can rely on high-reach, low-persistence channels.
\end{remark}

\begin{definition}[Channel Portfolio]
\label{def:portfolio}
\index{channel portfolio}
A \emph{channel portfolio} for belief $i$ is the set of channels through which funding is directed, together with the allocation across them: $\{(k, \funding{k})\}_{k \in K_i}$ where $K_i \subseteq \{1, \ldots, m\}$.
A diversified portfolio spreads investment across multiple channel types to exploit different reach-credibility-persistence profiles simultaneously.
\end{definition}


\section{Budgets as Resource Allocation Vectors}
\label{sec:ch4-budgets}

\begin{definition}[Doxonomic Budget]
\label{def:budget}
\index{budget}
A \emph{doxonomic budget} for belief $i$ is a vector $\fundingvec_i = (\funding{1}, \ldots, \funding{m}) \in \R^m_{\geq 0}$ specifying the resource allocation across $m$ channels directed at promoting belief $i$.
The total doxonomic expenditure on $i$ is $F_i = \|\fundingvec_i\|_1 = \sum_k \funding{k}$.
\end{definition}

The budget allocation problem is: given a total expenditure $F_i$, how should a funder allocate across channels to maximize $\belief{i}$?
This is a constrained optimization problem whose solution depends on the relative cost-effectiveness of each channel.

\begin{proposition}[Optimal Allocation]
\label{prop:optimal-allocation}
Under the first-order model~\eqref{eq:first-order} with concave cost-reach functions $c_k$, the optimal allocation satisfies the equimarginal condition:
\[
  \credibility{k} \cdot \relevance{k}{i} \cdot c_k'(\funding{k}) \cdot \persistence{k} = \lambda \quad \text{for all } k \text{ with } \funding{k} > 0
\]
where $\lambda$ is the Lagrange multiplier on the budget constraint $\sum_k \funding{k} = F_i$.
\end{proposition}

\begin{proof}
The funder maximizes $\belief{i} = \sum_k \funding{k} \cdot \credibility{k} \cdot \relevance{k}{i} \cdot c_k(\funding{k}) / \funding{k} \cdot \persistence{k}$ subject to $\sum_k \funding{k} = F_i$ and $\funding{k} \geq 0$.
Replacing $\reach{k}$ with $c_k(\funding{k})$ (the cost-to-reach function) and forming the Lagrangian $\mathcal{L} = \sum_k \credibility{k} \cdot \relevance{k}{i} \cdot c_k(\funding{k}) \cdot \persistence{k} - \lambda(\sum_k \funding{k} - F_i)$, the first-order condition for interior solutions is $\credibility{k} \cdot \relevance{k}{i} \cdot c_k'(\funding{k}) \cdot \persistence{k} = \lambda$.
Concavity of $c_k$ ensures the second-order conditions are satisfied.
\end{proof}

This is the doxonomic analogue of the equimarginal principle in economics: at the optimum, the marginal belief return per dollar is equalized across all active channels.
Channels where the marginal return exceeds $\lambda$ should receive more funding; channels where it falls below $\lambda$ should be cut.

\begin{example}[Budget Allocation: Think Tank vs.\ Advertising]
\label{ex:allocation}
A funder with $F = \$5\text{M}$ considers two channels:
\begin{itemize}[nosep]
  \item Think tank: $\credibility{1} = 0.75$, $\persistence{1} = 10$ years, $c_1(f) = 0.02\sqrt{f}$ (low reach, high credibility).
  \item Social media advertising: $\credibility{2} = 0.2$, $\persistence{2} = 0.1$ years, $c_2(f) = 0.5\sqrt{f}$ (high reach, low credibility).
\end{itemize}
The marginal return at funding level $f$ is $\credibility{k} \cdot \persistence{k} \cdot c_k'(f)$.
For the think tank: $0.75 \cdot 10 \cdot 0.01 / \sqrt{f} = 0.075/\sqrt{f}$.
For advertising: $0.2 \cdot 0.1 \cdot 0.25/\sqrt{f} = 0.005/\sqrt{f}$.
The think tank's marginal return is 15 times higher.
At the optimum, the funder should invest almost entirely in the think tank, a non-obvious result that reflects the enormous premium of credibility $\times$ persistence over raw reach.
\end{example}

\begin{remark}[Why Doxonomic Budgets Are Hard to Measure]
\label{rem:budget-opacity}
In practice, doxonomic budgets are difficult to estimate because:
\begin{itemize}[nosep]
  \item much funding flows through intermediaries (foundations, donor-advised funds) that obscure the ultimate source,
  \item institutional overhead blurs the distinction between doxonomic spending and general operations,
  \item some channels (education, entertainment) have doxonomic effects that are side effects of their primary function, not their explicit purpose,
  \item non-monetary resources (volunteer labor, social media sharing, cultural prestige) contribute to belief production but are not captured in financial accounting.
\end{itemize}
Chapter~\ref{ch:belief-flow-accounting} develops methods for addressing these measurement challenges.
\end{remark}


\section{Narratives}
\label{sec:ch4-narratives}

Narratives are the \emph{content} of doxonomic systems: what institutions produce, channels transmit, and audiences process.

\begin{definition}[Narrative]
\label{def:narrative}
\index{narrative}
A \emph{narrative} $n \in \narspace$ is a structured proposition or set of propositions equipped with:
\begin{enumerate}[nosep]
  \item a \emph{causal claim} (``$X$ because $Y$''): poverty exists because people don't work hard enough; markets crashed because of government interference; crime rises because of immigration,
  \item a \emph{normative framing} (``this is good/bad/natural/inevitable''): inequality is natural; growth is progress; austerity is responsible,
  \item a \emph{subject position} (who the audience is invited to identify with): the hardworking taxpayer, the innovative entrepreneur, the responsible citizen.
\end{enumerate}
\end{definition}

The same factual content can be embedded in radically different narratives.
``Unemployment is at 8\%'' becomes:
\begin{itemize}[nosep]
  \item $n_1$: ``The economy is recovering: unemployment has fallen from 10\% to 8\%.'' (Causal: recovery. Frame: progress. Subject: optimistic citizen.)
  \item $n_2$: ``Eight million people cannot find work, and the government has failed them.'' (Causal: policy failure. Frame: injustice. Subject: abandoned worker.)
  \item $n_3$: ``Welfare dependency keeps people from seeking employment.'' (Causal: moral failure. Frame: personal responsibility. Subject: taxpayer.)
\end{itemize}
Same statistic, three different doxonomic effects.
The narrative, not the fact, is the doxonomic product.

\begin{definition}[Narrative Distance]
\label{def:narrative-distance}
\index{narrative distance}
Given a vector-space representation of narratives (e.g., via document embeddings), the \emph{narrative distance} between two narratives $n_1, n_2 \in \narspace$ is
\[
  d(n_1, n_2) = 1 - \frac{\langle \bm{v}_{n_1}, \bm{v}_{n_2} \rangle}{\|\bm{v}_{n_1}\| \cdot \|\bm{v}_{n_2}\|}
\]
where $\bm{v}_n$ is the embedding vector of narrative $n$.
Large narrative distance between the dominant frame and an alternative frame indicates that the alternative is not merely a different opinion but a \emph{different way of seeing}, harder to reach by incremental argument, requiring more doxonomic investment to install.
\end{definition}

\begin{remark}[Narratives vs.\ Propositions]
\label{rem:narrative-proposition}
A narrative is more than a proposition.
The proposition ``inequality is increasing'' has a truth value.
The narrative ``inequality is increasing because the lazy are being rewarded'' has a truth value, a moral judgment, a causal attribution, and a subject position.
Doxonomic systems rarely produce naked propositions; they produce narratives, because narratives engage identity, emotion, and action in ways that propositions alone cannot.
This is why doxonomic production operates at multiple layers of the ontology simultaneously (Principle~\ref{prin:layered-ontology}).
\end{remark}


\section{Audiences and Prior Distributions}
\label{sec:ch4-audiences}

Audiences are not homogeneous.
A population consists of subgroups $\{G_1, \ldots, G_L\}$ with different prior belief distributions $\mu_{G_l}$, different media exposure patterns, different resistance profiles, and different identity commitments.
Effective doxonomic strategy requires targeting: knowing which audiences are persuadable and through which channels.

\begin{definition}[Audience Heterogeneity]
\label{def:audience-heterogeneity}
\index{audience}
The population-level belief $\belief{i}$ is a mixture:
\[
  \belief{i} = \sum_{l=1}^{L} w_l \cdot \belief{i}^{(l)}
\]
where $w_l$ is the population share of group $l$ and $\belief{i}^{(l)}$ is the average credence in group $l$.
A doxonomic system may target some groups more than others.
\end{definition}

\begin{definition}[Audience Resistance]
\label{def:resistance}
\index{audience resistance}
The \emph{resistance} of audience group $G_l$ to narrative $n$ is a function $\rho_l(n) \in [0,1]$ reflecting the degree to which the group's prior beliefs, identity commitments, lived experience, and critical resources reduce the narrative's doxonomic effect.
The effective belief change in group $l$ from exposure to $n$ through channel $k$ is
\[
  \Delta\belief{i}^{(l)} = (1 - \rho_l(n)) \cdot \credibility{k} \cdot \relevance{k}{i} \cdot \persistence{k}
\]
High resistance ($\rho_l \to 1$) means the group is effectively immune to the narrative; low resistance ($\rho_l \to 0$) means full susceptibility.
\end{definition}

Resistance is not random.
It is shaped by:
\begin{itemize}[nosep]
  \item \textbf{Lived experience}: a worker who has been laid off despite ``working hard'' is more resistant to meritocratic narratives than a comfortably employed professional.
  \item \textbf{Counter-institutions}: groups with access to alternative information sources (unions, community organizations, alternative media) have higher resistance.
  \item \textbf{Identity}: narratives that threaten core identity commitments provoke backlash rather than persuasion \autocite{kahan2012}.
  \item \textbf{Prior exposure}: groups that have been ``inoculated'' against a narrative (Chapter~\ref{ch:counter-doxonomic}) show elevated resistance to subsequent exposure.
\end{itemize}

\begin{example}[Differential Resistance to Meritocratic Narratives]
\label{ex:resistance}
Survey data consistently shows that belief in meritocracy is negatively correlated with direct experience of structural barriers.
Black Americans score lower on meritocratic belief scales than white Americans; women in male-dominated workplaces are more skeptical of ``effort determines success'' than men in the same workplaces; first-generation college students are less convinced of equal opportunity than their legacy peers.
In doxonomic terms, lived experience of structural disadvantage raises $\rho_l$ for meritocratic narratives, a form of experiential resistance that no amount of doxonomic investment can fully overcome.
\end{example}


\section{Feedback Loops}
\label{sec:ch4-feedback}

The belief production chain is not unidirectional.
Beliefs, once internalized, alter behavior, which alters institutions, which alters funding.
This feedback structure is what makes doxonomic systems self-sustaining, and what makes them so difficult to disrupt.

\begin{definition}[Doxonomic Feedback]
\label{def:feedback}
\index{feedback loop}
A \emph{doxonomic feedback loop} is a dynamical coupling in which the belief state $\beliefvec(t)$ influences the funding vector $\fundingvec(t+1)$ through behavioral and institutional responses:
\[
  \fundingvec(t+1) = g(\beliefvec(t), \fundingvec(t))
\]
for some function $g$.
The system is \emph{self-reinforcing} if $\partial g_k / \partial \belief{i} > 0$ for the belief $i$ that channel $k$ promotes.
The system is \emph{self-undermining} if $\partial g_k / \partial \belief{i} < 0$ (success reduces future investment).
\end{definition}

\begin{example}[Self-Reinforcing: The Tax Cut Cycle]
\label{ex:tax-cycle}
\begin{enumerate}[nosep]
  \item Think tanks promote the belief that tax cuts stimulate growth ($\belief{T}$ increases).
  \item Population supports tax-cutting politicians (behavioral consequence).
  \item Tax cuts are enacted; wealthy beneficiaries gain additional disposable income.
  \item Beneficiaries donate to the same think tanks (funding increases).
  \item Think tanks produce more tax-cut advocacy ($\belief{T}$ increases further).
\end{enumerate}
This is a self-reinforcing loop: $\partial g_k / \partial \belief{T} > 0$.
Each iteration strengthens the cycle.
\end{example}

\begin{example}[Self-Undermining: The Backlash Effect]
\label{ex:backlash}
\begin{enumerate}[nosep]
  \item A corporation promotes the belief that its product is environmentally friendly (``greenwashing'').
  \item The population initially accepts the claim ($\belief{G}$ increases).
  \item Investigative journalists or activists expose the deception.
  \item Public trust in the corporation's claims collapses ($\credibility{k}$ drops).
  \item Further spending on the same message becomes counterproductive.
\end{enumerate}
This is a self-undermining loop: success (increased scrutiny) reduces future effectiveness.
\end{example}

Self-reinforcing feedback is the mechanism by which temporary doxonomic investments become permanent features of a society's common sense.
Self-undermining feedback is the mechanism by which doxonomic systems can fail, a topic developed in Chapter~\ref{ch:how-belief-regimes-fracture}.

\begin{definition}[Doxonomic Stability]
\label{def:stability}
\index{doxonomic stability}
A doxonomic system is \emph{locally stable} at belief state $\beliefvec^*$ if small perturbations (temporary counter-narratives, one-time contrary evidence) are absorbed by the feedback structure and the system returns to $\beliefvec^*$.
Formally, $\beliefvec^*$ is a stable fixed point of the composed dynamical system:
\[
  \beliefvec^* = \Psi(\Phi(g(\beliefvec^*, \fundingvec^*)), \beliefvec^*)
\]
where $\Phi$ is the production function, $\Psi$ is the internalization function, and $g$ is the feedback function.
\end{definition}


\section{Assembling the Units: A Toy Model}
\label{sec:ch4-toy}

To see how the units work together, consider a minimal doxonomic system with one funder, one institution, one channel, and a homogeneous audience.

\begin{model}[Minimal Doxonomic System]
\label{mod:minimal}
\begin{itemize}[nosep]
  \item Funder: budget $F$, objective $\max \belief{}$.
  \item Institution: processing function $\iota(F) = (n, \credibility{})$ where $\credibility{} = 0.8$.
  \item Channel: reach $c(F) = 1 - e^{-F/F_0}$ (saturating), persistence $T = 5$ years.
  \item Audience: $N = 1000$ agents, prior $\pi_j^{(0)} = 0.3$, Bayesian updating.
  \item Feedback: $F(t+1) = F_0 + \beta \belief{}(t)$ (self-reinforcing with rate $\beta$).
\end{itemize}
The dynamics are:
\[
  \belief{}(t+1) = \belief{}(t) + \alpha \cdot \credibility{} \cdot c(F(t)) \cdot T \cdot (1 - \belief{}(t)) - \delta \cdot \belief{}(t)
\]
where $\alpha$ is the persuasion rate and $\delta$ is natural decay.

This system has at most two stable fixed points: a low-belief equilibrium (the narrative fails to take hold) and a high-belief equilibrium (common-sense capture).
The basin of attraction depends on the initial funding $F_0$ and the feedback strength $\beta$.
\end{model}

This toy model is developed into a full dynamical framework in Chapters~\ref{ch:formal-models} and~\ref{ch:belief-production-chain}.


\section{Notes and References}
\label{sec:ch4-notes}

Agent-based models of belief formation draw on \textcite{jackson2008} and \textcite{granovetter1978}.
The institutional processing framework extends \textcite{north1990} on institutions as ``rules of the game.''
On institutional autonomy and the relative independence of cultural production, see \textcite{bourdieu1984}.
Optimal allocation of persuasion resources connects to \textcite{downs1957} and \textcite{tullock1967}.
Audience heterogeneity models are treated in \textcite{manski1993}; on identity-protective cognition, see \textcite{kahan2012}.
Feedback loops in social systems are central to \textcite{schelling1978}.
The climate-doubt example draws on \textcite{oreskes2010}.
On the economics of attention as a scarce resource, see \textcite{kahneman2011}.
Credibility laundering through institutional intermediaries is documented in \textcite{mayer2016}.
Narrative as a structured object with causal, normative, and identity components draws on framing theory; see \textcite{stanley2015} for a philosophical treatment.


\section{Exercises}

\begin{exercise}
Classify the following as agents, institutions, channels, or narratives: (a)~the Koch Foundation, (b)~the \textit{New York Times}, (c)~Twitter/X, (d)~``free markets create prosperity,'' (e)~a university economics department, (f)~a YouTube algorithm, (g)~a public school curriculum, (h)~``hard work is its own reward.''
Note cases where an entity plays multiple roles.
\end{exercise}

\begin{exercise}
\label{ex:ch4-allocation}
Compute the optimal budget allocation for a funder with $F_i = \$10\text{M}$ across two channels with $\credibility{1} = 0.9$, $\credibility{2} = 0.5$, and cost-reach functions $c_1(f) = \sqrt{f/10^6}$, $c_2(f) = 2\sqrt{f/10^6}$.
Assume $\relevance{k}{i} = \persistence{k} = 1$ for both.
How does the allocation change if $\persistence{1} = 20$ and $\persistence{2} = 0.5$?
\end{exercise}

\begin{exercise}
Model a self-reinforcing feedback loop in which belief in meritocracy increases tolerance for inequality, which increases political donations from wealthy actors, which increases funding for meritocratic narratives.
Write the system as a pair of coupled difference equations and identify the conditions for a stable high-meritocracy equilibrium.
\end{exercise}

\begin{exercise}
Consider a population with two groups: $G_1$ (30\% of population, high media exposure, low prior credence in $p$, low resistance $\rho_1 = 0.2$) and $G_2$ (70\%, low media exposure, high prior credence in $p$, high resistance $\rho_2 = 0.8$).
Under what conditions does a doxonomic investment targeting $G_1$ change the population-level $\belief{i}$ more than one targeting $G_2$?
\end{exercise}

\begin{exercise}
For the minimal doxonomic system (Model~\ref{mod:minimal}), find the fixed points of the dynamics as a function of $\alpha$, $\delta$, $\beta$, $F_0$, and $\credibility{}$.
Under what conditions does only one fixed point exist (inevitable capture or inevitable failure)?
Under what conditions do two stable fixed points coexist?
\end{exercise}

\begin{exercise}
Design a doxonomic system from scratch: choose a belief $p$, specify funders, institutions, channels, narratives, target audiences, and feedback structure.
Estimate plausible values for $\credibility{k}$, $\reach{k}$, $\persistence{k}$, and $\rho_l$.
Compute the expected $\belief{i}$ using the first-order model.
Then identify the weakest link in your system, the point where counter-doxonomic intervention would be most effective.
\end{exercise}

\begin{exercise}
The chapter notes that entertainment operates at layers 3--5 of the ontology.
Choose a popular television series or film franchise and analyze it as a doxonomic channel.
What anthropological regime does it normalize?
What subject positions does it offer?
Estimate its $\reach{}$, $\credibility{}$ (noting that entertainment credibility works differently from news credibility), and $\persistence{}$.
\end{exercise}
